% !PRINTED 556
\weberpage{556}

\weq{6}{\varphi(r) = (b_1 - b_0)\,r + (b_2 - b_0)\,r^2 + \cdots + (b_{n-2} - b_0)\,r^{n-2} - b_0 r^{n-1};}

or, since the powers $r, r^2, \ldots r^{n-1}$, apart from their order, agree with $r, r_1, r_2 \ldots r_{n-2}$, into the form

\weq{7}{\varphi(r) = a\,r + a_1 r_1 + a_2 r_2 + \cdots + a_{n-2} r_{n-2},}

in which the coefficients $b$ and $a$ are rational numbers.

\gesperrt{I. Such a function $\varphi(r)$ can be equal to zero only when all its coefficients $a, a_1 \ldots a_{n-2}$ vanish;}

for the expression (5) shows that $\varphi(r)$ cannot vanish otherwise than when the $b_0, b_1 \ldots b_{n-2}$ are zero. But if these are zero, then (6) shows that the coefficients $a$ must also all be zero. Hence $\varphi(r)$ can be brought into the form (7) in only one way.

If we make the stipulation that $a_h$ and $r_h$ shall not change when the index $h$ increases by a multiple of $n - 1$, we can also set

\weq{8}{\varphi(r) = \sum^h a_h r_h}

and let $h$ run through a complete system of residues with respect to the modulus $n - 1$.

We now have to prove that the group of the equation $X = 0$ is no other than the period of the cyclic permutation
\[
\pi = (r, r_1, r_2 \ldots r_{n-2}).
\]

This period or cyclic group we denote by

\weq{9}{C = 1, \pi, \pi^2 \ldots \pi^{n-2}.}

The equation $X = 0$ is in any case a \gesperrt{normal equation}, because it is irreducible, and because all its roots are, by (3), rationally expressible in terms of one among them; it is therefore its own \gesperrt{Galois} resolvent (\S.\,145). Its substitutions are

\weq{10}{(r, r), (r, r_1), (r, r_2) \ldots (r, r_{n-2}),}

and by (3) we have
\[
(r, r_h) = (r_k, r_{k+h}).
\]

From this it follows that $(r, r_1)$ produces among the roots (4) the cyclic permutation $\pi$, and that therefore the substitution group
%CONT>>
